\kchapter{Speaking Kubernetes: kubectl and YAML}{api}

\begin{objectives}
\begin{itemize}
  \item just enough YAML to read and write any manifest;
  \item the four fields every Kubernetes object has;
  \item imperative commands vs.\ declarative \kc{kubectl apply};
  \item the twenty \kc{kubectl} commands you will use every day;
  \item how to discover fields yourself with \kc{kubectl explain}.
\end{itemize}
\end{objectives}

\section{YAML in ten minutes}

YAML is a human-friendly way to write structured data. Kubernetes accepts JSON
too, but everyone writes YAML. There are only a few rules.

\begin{yamlcode}{yaml-basics.yaml}
# A comment
name: web                 # key: value (a "map" entry)
replicas: 3               # numbers are not quoted
enabled: true             # booleans
version: "1.0"            # quote strings that look like numbers

labels:                   # a nested map (indent with SPACES, usually 2)
  app: web
  tier: frontend

ports:                    # a list: each item starts with "- "
  - 80
  - 443

containers:               # a list of maps
  - name: nginx
    image: nginx:1.27
  - name: sidecar
    image: busybox:1.36

script: |                 # multi-line string, newlines kept
  echo hello
  echo world
---
# "---" separates multiple documents in one file
kind: Second
\end{yamlcode}

\begin{warningbox}
\textbf{Never use tabs} in YAML, and be careful with indentation: it defines
structure. A key indented by one extra space moves into a different object, and
the error message will be confusing. Use an editor with a YAML extension.
\end{warningbox}

\section{Anatomy of a manifest}

Every Kubernetes object has the same skeleton:

\begin{yamlcode}{pod.yaml}
apiVersion: v1            # 1. which API group/version defines this kind
kind: Pod                 # 2. what type of object
metadata:                 # 3. identity: name, namespace, labels, annotations
  name: hello
  namespace: default
  labels:
    app: hello
spec:                     # 4. desired state (different for every kind)
  containers:
    - name: hello
      image: nginx:1.27
      ports:
        - containerPort: 80
\end{yamlcode}

\begin{center}
\ktablestart
\begin{tabularx}{\linewidth}{P{2.6cm} L}
\kheader \kth{Field} & \kth{Explanation} \\
\kc{apiVersion} & Core objects use \kc{v1}. Others are \kc{group/version},
                  e.g.\ \kc{apps/v1}, \kc{batch/v1}, \kc{networking.k8s.io/v1}. \\
\kc{kind} & The object type: Pod, Deployment, Service, ConfigMap\ldots \\
\kc{metadata} & \kc{name} must be unique per kind per namespace. Labels and
                annotations live here. \\
\kc{spec} & What you want. Kubernetes adds a \kc{status} section later. \\
\end{tabularx}
\end{center}

\section{Discovering the API yourself}

You never need to memorise fields. Ask the cluster:

\begin{bashcode}{terminal}
kubectl api-resources                    # every kind, its short name, and API group
kubectl explain pod                      # docs for a kind
kubectl explain pod.spec.containers      # drill into fields
kubectl explain deployment.spec.strategy --recursive
\end{bashcode}

Short names save typing: \kc{po} (pods), \kc{deploy}, \kc{svc}, \kc{ns},
\kc{cm} (configmaps), \kc{pvc}, \kc{sts}, \kc{ds}, \kc{ing}, \kc{hpa}, \kc{sa}.

\section{Imperative vs.\ declarative}

\begin{bashcode}{terminal}
# Imperative: tell Kubernetes what to DO (fine for quick experiments)
kubectl run hello --image=nginx:1.27
kubectl create deployment web --image=nginx:1.27 --replicas=3
kubectl scale deployment web --replicas=5
kubectl delete deployment web

# Declarative: tell Kubernetes what should EXIST (use this for real work)
kubectl apply -f pod.yaml          # create or update to match the file
kubectl apply -f ./manifests/      # a whole folder
kubectl diff  -f pod.yaml          # preview what would change
kubectl delete -f pod.yaml
\end{bashcode}

\begin{keyidea}
\textbf{Files in Git are the source of truth; \kc{kubectl apply} makes the
cluster match them.} You can re-run \kc{apply} a hundred times and nothing
changes unless the file changes. This property (idempotence) is what makes
GitOps possible later.
\end{keyidea}

\begin{protip}
Let \kc{kubectl} write YAML for you, then edit it:
\begin{bashcode}{terminal}
kubectl create deployment web --image=nginx:1.27 --dry-run=client -o yaml > web.yaml
\end{bashcode}
\end{protip}

\section{The everyday command set}

\begin{bashcode}{terminal}
# --- Looking around -------------------------------------------------------
kubectl get pods                         # list in the current namespace
kubectl get pods -A                      # all namespaces
kubectl get pods -o wide                 # + node and IP
kubectl get pod hello -o yaml            # full object, including status
kubectl get deploy,svc,pods -l app=web   # several kinds, filtered by label
kubectl get pods -w                      # watch for changes
kubectl describe pod hello               # human summary + EVENTS (read these!)
kubectl get events --sort-by=.lastTimestamp

# --- Debugging ------------------------------------------------------------
kubectl logs hello                       # container stdout/stderr
kubectl logs hello -c sidecar -f         # a specific container, follow
kubectl logs deploy/web --previous       # logs of the crashed previous run
kubectl exec -it hello -- sh             # open a shell inside the container
kubectl port-forward pod/hello 8080:80   # localhost:8080 -> pod port 80
kubectl debug -it hello --image=busybox:1.36 --target=hello   # ephemeral debug container
kubectl top pods                         # CPU/memory (needs metrics-server)

# --- Changing things ------------------------------------------------------
kubectl apply -f file.yaml
kubectl edit deployment web              # edit live object in your $EDITOR
kubectl set image deployment/web nginx=nginx:1.28
kubectl label pod hello env=dev
kubectl annotate pod hello owner=team-a
kubectl delete pod hello
\end{bashcode}

\subsection{Output formats}

\begin{bashcode}{terminal}
kubectl get pods -o name
kubectl get pods -o jsonpath='{.items[*].metadata.name}'
kubectl get pods -o custom-columns=NAME:.metadata.name,NODE:.spec.nodeName
kubectl get pod hello -o json | jq '.status.podIP'
\end{bashcode}

\begin{tryit}
\begin{enumerate}
  \item Save \kc{pod.yaml} above and run \kc{kubectl apply -f pod.yaml}.
  \item Run \kc{kubectl get pod hello -o yaml} and find the \kc{status} section
        Kubernetes added. What is \kc{status.podIP}?
  \item Port-forward with \kc{kubectl port-forward pod/hello 8080:80} and open
        \url{http://localhost:8080}.
  \item Change the image to \kc{nginx:1.28} in the file and apply again. What
        happens? (Hint: some Pod fields are immutable --- read the error.)
  \item Clean up with \kc{kubectl delete -f pod.yaml}.
\end{enumerate}
\end{tryit}

\begin{recap}
\begin{itemize}
  \item Manifests = \kc{apiVersion} + \kc{kind} + \kc{metadata} + \kc{spec}.
  \item Use \kc{kubectl explain} and \kc{api-resources} instead of memorising.
  \item Prefer declarative \kc{apply} with files in Git.
  \item \kc{describe} and \kc{get events} are your first debugging tools.
\end{itemize}
\end{recap}
